\section{Scope of the raw fourth-moment reduction}
\label{upq:details}
Write $c_{a,b}=\mu(a)\mu(b)\nu(ab)W(Nab/D)$ on the exact factor support
in the core. Before absolute Cauchy, the normalized energy is
\[
 \frac1D\sum_{\substack{(a,b),(c,d)\ {\rm on\ that\ support}}}
 c_{a,b}\overline{c_{c,d}}
 \sum_{u\in B}x_a(u)x_b(u)\overline{x_c(u)x_d(u)}.
\]
Squarefreeness of $ab$ and $cd$ forces each positive and negative
prime multiplicity to be at most one, proving $e_p\in\{0,\pm1\}$.
For the mixed term against the actual $\mu$ polynomial the exponent is
$\mathbf1_{p\mid a}+\mathbf1_{p\mid b}-\mathbf1_{p\mid n}$,
with the same conclusion and the same full union mask.
The common possibly complex $W$, $\nu$, hyperbolic support and all raw
zeros stay in these identities.

In contrast, unrestricted four-column memberships $(1,1,0,0)$ give
$e_p=2$, a cubic component; $(1,1,1,0)$ give $e_p=1$ although both
product supports contain $p$. Neither the parity symmetric difference
of four supports nor the symmetric difference of two nonsquarefree
product supports gives the general conductor.
For the interface in the core, its local character order is
\[
 \begin{array}{c|cccc}
 j&1,5&2,4&3&6\\ \hline
 e_p=\pm1&6&3&2&1\\
 e_p=\pm2&3&3&1&1
 \end{array}
\]
Nontrivial entries have tame conductor exponent one; trivial entries
remain deleted at $p\mid v$. Multiplying the interface over all six
row factors, with the fixed unit/$S$ phase, preserves every original
zero, including those at sixth-power factors. All inter-factor,
squarefree, ray and capacity restrictions remain external restrictions
on the same row sum; this identity asserts no independence.

At a residue field of size $Q$, let $\chi$ have order six and $\rho$
be nontrivial additive. The zero-extended complete sum
$G_k(t)=\sum_{z\ne0}\chi(z)^k\rho(tz)$ satisfies
\[
 \begin{array}{c|cc}
 &t=0&t\ne0\\ \hline
 k\not\equiv0\ (6)&0&|G_k(t)|=\sqrt Q\\
 k\equiv0\ (6)&Q-1&-1.
 \end{array}
\]
For $t\ne0$, substitution gives
$G_k(t)=\overline{\chi(t)^k}G_k(1)$; for nontrivial $k$,
expanding the norm square gives
$(Q-1)-\sum_{a\ne0,1}\chi(a)^k=Q$.
The proof applies at split and inert primes. By CRT the full period-$L$
transform vanishes if any local frequency on $Q_j$ is zero; otherwise
its magnitude is
\[
 (NQ_j)^{1/2}
 \prod_{p\mid Z_j}
 \left|-1+Np\,\mathbf1_{\text{local frequency at }p\text{ is }0}\right|.
\]
Canceled common primes thus supply Ramanujan factors.

Exact deletion, with all remaining restrictions kept in $F$, gives
\[
 \sum_r\theta(r)F(r)\mathbf1_{(r,Z)=1}
 =\sum_{d\mid Z}\mu(d)\theta(d)\sum_z\theta(z)F(dz).
\]
For a squarefree pivot let $E$ be the radical of its fixed excluded
primes, including $Z$ and the relevant other row factors, with
conductor primes omitted since $\theta$ already vanishes there.
The additional squarefree identity is
\[
 \mu(r)^2\mathbf1_{(r,E)=1}
 =\sum_{\substack{h^2\mid r\\(h,E)=1}}\mu(h)
       \sum_{\substack{d\mid E\\d\mid r}}\mu(d).
\]
Since $(d,h)=1$, its lifted row sum has terms
$\mu(d)\mu(h)\theta(d)\theta(h)^2\sum_z\theta(z)F(dh^2z)$.
The scaled $F$ retains the exact window, primary/ray restrictions and
other masks. Divisor triangles cost at most $(HD)^\epsilon$;
they do not restore a forbidden nonzero.

The previous fixed-radial-profile high-gcd proof
(\path{research/short-row-high-gcd.md}, independently reviewed in
\path{planning/reviews/book/short-row-formal.md})
already bounds its off-diagonal tail by
$O_{\Phi,\epsilon}(D^{2+\epsilon}/G_0^2)$ for $Nh\ge G_0$.
For annular $m=hr,n=ht$, $Nq\le Y$ implies $Nh\ge D/\sqrt Y$,
so the old smooth-scope bound is $O_{\Phi,\epsilon}(D^\epsilon Y)$.
The present elementary count $O(D\sqrt Y)$, factor multiplicities
$r_U(m)r_U(n)\ll_\epsilon D^\epsilon$, and $\#B_A=O(R_0)$ give only
$O(\|W\|_\infty^2D^\epsilon R_0\sqrt Y)$ on a literal sharp block.
This is a weaker conductor-form recovery in its own bookkeeping scope,
not a new high-gcd theorem; the smooth estimate is not silently
transferred to a selected sharp off-diagonal piece.
The full $A_3$ face and neighborhood slack are also inherited.

Finally, primitive smooth Poisson completion has bound
$O_V(\min(A_j,\sqrt{Nq}))$ on the original support.
For $Nq>P^2$ and $1\le j\le5$, every scaled length
$A_j/(Nd\,Nh^2)$ is below $\sqrt{Nq}$, so even this favorable trial
takes its trivial length term. Its constant uses fixed smooth
seminorms and need not stay bounded when a sharp boundary is approached.
In the admissible subsector $(m,n)=1$, $Nq=Nm\,Nn\asymp D^2$.
A complete zero mean does not control the nonzero frequencies of the
short original masked row window. The unresolved target is their
joint signed sum with weights $r_U(m)r_U(n)\lambda_m\overline{\lambda_n}$
and $r_U(m)\lambda_m\overline{\lambda_n}$, before per-conductor absolute
estimates. No new global analytic conclusion is claimed.
